\section{引言}
\label{sec:introduction}
本地论文编辑器不仅需要正确显示文本，还需要协调多文件输入、编译中间文件、参考文献辅助程序和 PDF 预览。一份仅有标题和一段文字的文档，难以暴露交叉引用未收敛、图片未被复制到快照或参考文献未生成等问题。因此，测试输入应具有接近论文的结构，同时使每一个数值和引用都可以追踪。

本样例选择传感器异常检测作为叙述载体。周期性温度序列具有直观的物理类比，阈值方法易于表达为公式，也能自然产生曲线图、混淆矩阵及指标表。这里的“温度”只是测试单位，不来源于真实传感器。数据集编号为 LOL-SYN-120，生成约定见内部测试说明\cite{lolData2026}。

本文围绕三个问题组织实验：第一，固定阈值与残差阈值在含周期项的序列上呈现何种差异；第二，阈值降低是否伴随误报与漏报的变化；第三，数据、表格、图片和引用能否在离线环境中组成可重复编译的论文。前两个问题只在本合成序列内讨论，第三个问题才是项目的主要验收目标。

\subsection{相关工作与引用边界}
为了避免将虚构文献包装为真实学术来源，本文仅引用三个本地内部测试记录。数据协议\cite{lolData2026}定义生成公式与标签；指标说明\cite{lolMetrics2026}固定混淆矩阵方向及舍入规则；复现指南\cite{lolReproduce2026}描述脚本、文件结构及编译顺序。这些条目用于验证 BibTeX 引用链路，不构成外部文献综述，也不能支持现实世界中的算法有效性结论。

\subsection{组织结构}
第\ref{sec:method}节定义数据和检测规则，第\ref{sec:experiment}节展示数值与图像，第\ref{sec:discussion}节讨论限制，附录提供全部观测记录。跨文件章节引用、表格引用及公式引用均由 LaTeX 自动编号，禁止手工写死页码和序号。
